% !TeX program = xelatex
% DCGI v0.3 — 理由档案
% Compile with: latexmk -xelatex -interaction=nonstopmode -halt-on-error DCGI_v0.3.tex
% This document is self-contained. No external fonts or images are bundled.
\documentclass[UTF8,zihao=-4,fontset=none,oneside,openany]{ctexbook}
\usepackage[paperwidth=176mm,paperheight=250mm,left=22mm,right=21mm,top=24mm,bottom=23mm,headheight=15pt,headsep=8mm,footskip=11mm]{geometry}
\usepackage{fontspec}
\IfFontExistsTF{Noto Serif CJK SC}{%
  \setCJKmainfont{Noto Serif CJK SC}[BoldFont={Noto Serif CJK SC Bold},ItalicFont={Noto Serif CJK SC}]
  \setCJKsansfont{Noto Sans CJK SC}[BoldFont={Noto Sans CJK SC Bold}]
}{%
  \setCJKmainfont{FandolSong-Regular}[Extension=.otf,BoldFont=FandolSong-Bold,ItalicFont=FandolKai-Regular]
  \setCJKsansfont{FandolHei-Regular}[Extension=.otf,BoldFont=FandolHei-Bold]
}
\setmainfont{lmroman10-regular.otf}[BoldFont=lmroman10-bold.otf,ItalicFont=lmroman10-italic.otf,BoldItalicFont=lmroman10-bolditalic.otf]
\setsansfont{lmsans10-regular.otf}[BoldFont=lmsans10-bold.otf,ItalicFont=lmsans10-oblique.otf,Scale=MatchLowercase]
\setmonofont{lmmono10-regular.otf}[Scale=MatchLowercase]
\usepackage{amsmath,amsthm}
\usepackage{unicode-math}
\setmathfont{latinmodern-math.otf}
\usepackage{microtype}
\usepackage{setspace}
\setstretch{1.20}
\setlength{\parindent}{2em}
\setlength{\parskip}{2.5pt plus 0.4pt}
\setlength{\emergencystretch}{2em}
\clubpenalty=10000
\widowpenalty=10000
\displaywidowpenalty=10000
\raggedbottom
\usepackage{xcolor}
\definecolor{ink}{HTML}{263B4B}
\definecolor{muted}{HTML}{61717B}
\definecolor{gold}{HTML}{A68452}
\definecolor{wash}{HTML}{F3F5F6}
\usepackage{booktabs,longtable,tabularx,array}
\usepackage{enumitem}
\setlist{nosep,leftmargin=2em,topsep=4pt}
\usepackage{needspace}
\usepackage[most]{tcolorbox}
\newtcolorbox{principle}{enhanced,breakable,colback=wash,colframe=ink,boxrule=0pt,borderline west={1.5pt}{0pt}{ink},arc=0mm,left=11pt,right=10pt,top=9pt,bottom=9pt,before skip=11pt,after skip=12pt}
\newtcolorbox{casequote}{enhanced,breakable,colback=white,colframe=gold,boxrule=0pt,borderline west={1.2pt}{0pt}{gold},arc=0mm,left=12pt,right=7pt,top=5pt,bottom=5pt,before skip=9pt,after skip=9pt}
\usepackage{fancyhdr}
\pagestyle{fancy}
\fancyhf{}
\fancyhead[L]{\small\sffamily\color{muted}DCGI\quad 闭依开生法}
\fancyhead[R]{\small\sffamily\color{muted}\nouppercase{\leftmark}}
\fancyfoot[C]{\small\color{muted}\thepage}
\renewcommand{\headrulewidth}{0.3pt}
\renewcommand{\footrulewidth}{0pt}
\renewcommand{\chaptermark}[1]{\markboth{#1}{}}
\fancypagestyle{plain}{\fancyhf{}\fancyfoot[C]{\small\color{muted}\thepage}\renewcommand{\headrulewidth}{0pt}}
\ctexset{
  chapter={name={第,编},number=\chinese{chapter},format=\Huge\sffamily\bfseries\color{ink},nameformat=\Large\sffamily\color{muted},aftername=\par\vskip 10pt,beforeskip=15pt,afterskip=27pt},
  section={format=\Large\sffamily\bfseries\color{ink},beforeskip=22pt,afterskip=9pt},
  subsection={format=\normalsize\sffamily\bfseries\color{ink},beforeskip=16pt,afterskip=6pt}
}
\setcounter{secnumdepth}{0}
\setcounter{tocdepth}{1}
\usepackage[numbers,sort&compress]{natbib}
\usepackage{xurl}
\usepackage[unicode,bookmarksnumbered=true,bookmarksopen=true,bookmarksopenlevel=0,colorlinks=true,linkcolor=ink,citecolor=ink,urlcolor=ink]{hyperref}
\hypersetup{pdftitle={闭依开生法（DCGI）v0.3：思想前史、判例与规则资格档案},pdfauthor={Zheng Kuang},pdfsubject={Dependency Closure–Generative Independence Method},pdfkeywords={DCGI,闭依开生法,认知体系,路径原创,Return,just infiniteness,orbit equivalence},pdfcreator={XeLaTeX}}
\usepackage{bookmark}
\newcounter{dcgiarticle}
\newcommand{\Article}[1]{\Needspace{6\baselineskip}\refstepcounter{dcgiarticle}\subsection*{第\chinese{dcgiarticle}条\quad #1}\label{article:\arabic{dcgiarticle}}\pdfbookmark[2]{第\chinese{dcgiarticle}条　#1}{article-\arabic{dcgiarticle}}}
\newcommand{\lead}[1]{\par\noindent\textbf{#1}\par\nobreak\vspace{3pt}}
\newcommand{\rulebasis}[2]{\par\smallskip\noindent{\small\sffamily\color{muted}\textbf{来源：}#1\quad\textbf{压力测试状态：}#2}\par}
\newcommand{\claimstatus}[1]{\par\smallskip\noindent{\small\sffamily\color{muted}\textbf{当前状态：}#1}\par}
\newcommand{\Id}{\mathrm{Id}}
\newcommand{\AF}{\mathrm{AF}}
\newcommand{\Rist}{\operatorname{Rist}}
\newcommand{\rist}{\operatorname{rist}}
\newcommand{\Cof}{\operatorname{Cof}}
\newcommand{\IC}{\operatorname{IC}}
\newcommand{\BB}{\mathsf B}
\newcommand{\GG}{\mathfrak G}
\newcommand{\TT}{\mathfrak T_{\BB}}
\newcommand{\FF}{\mathsf F}
\newcommand{\standalonechapter}[1]{\chapter*{#1}\addcontentsline{toc}{chapter}{#1}\markboth{#1}{}}
\renewcommand{\bibname}{参考文献与阅读线索}
\begin{document}
\hypersetup{pageanchor=false}
\begin{titlepage}
\thispagestyle{empty}
\color{ink}
\vspace*{13mm}
\noindent{\sffamily\fontsize{12}{16}\selectfont DCGI\quad /\quad v0.3}\par
\vspace{8mm}
\noindent{\color{gold}\rule{25mm}{1.2pt}}\par
\vspace{14mm}
\noindent{\sffamily\bfseries\fontsize{36}{46}\selectfont 闭依开生法}\par
\vspace{7mm}
\noindent{\fontsize{15}{22}\selectfont 理解已有知识，释放新的可能}\par
\vspace{12mm}
\noindent{\sffamily\fontsize{11}{17}\selectfont Dependency Closure\par\noindent Generative Independence Method}\par
\vfill
\noindent{\fontsize{12}{21}\selectfont 思想前史\quad 宪法\quad 判例\quad 条款\quad AI 执行模板}\par
\vspace{7mm}
\noindent{\color{gold}\rule{\textwidth}{0.5pt}}\par
\vspace{6mm}
\noindent 匡正\quad Zheng Kuang\par
\vspace{2mm}
\noindent{\small\color{muted}2026 年 9 月\quad 理由档案}
\vspace*{5mm}
\end{titlepage}
\frontmatter
\hypersetup{pageanchor=true}
\standalonechapter{编写说明}
这份文稿保存一套在具体研究中逐渐形成、又被具体研究反过来修正的认知方法。它的名字是“闭依开生法”，英文简称 DCGI。v0.3 将可读版本明确定位为\textbf{理由档案}：它不负责替研究者逐步执行任务，而负责说明规则为什么存在、从哪里来、经受过什么压力，以及当前具有怎样的资格。执行层另以短模板派生，不与理由档案平行生长。

文中的“宪法”是方法论的比喻，指规定优先次序与适用边界的最高原则，不涉及法律制度。它之所以置于条款之前，是因为一套方法若只说明“可以怎样做”，却不说明“什么始终优先”，便很容易在执行中忘记自己的目的。

全文根据作者与 AI 的研究讨论重述。DCGI 并非从两次数学压力测试中突然出现。在此之前，彭国欣与 AI“钟上的降雪”围绕“路径原创”“问题先于概念”“语言—问题适配”“投影—路径—返回”“Monogenesis 与 Return”“去中心化后的重新整合”等主题形成了一组连续文章。它们提供了这套方法的思想前史：理解不是把现实迅速装入现成概念，而是在有限位置中生成自己的路径；语言应服务于问题而非取得垄断权；真正的 return 不是 rewind，而是把已经发生的差异与来路一起带回。本文将这些文章作为生成史材料，而不是用它们为后来的数学结论背书。\par

两次压力测试的先后顺序、关键提问、误判与纠正均予保留，叙述按问题的逻辑展开，而不是逐字复制聊天。第一次以 just infiniteness 为对象，第二次围绕 orbit equivalence 展开。它们不是用来装饰方法论的成功故事，而是拥有修改方法论的权力的案例。

为让不同背景的读者进入讨论，数学部分补入了必要的定义、短证明与文献线索。文献中的数学结果、讨论提出的候选路线，以及本书所作的方法论总结，各有自己的位置。\textbf{判例中的结构说明不替代原论文的完整假设与证明，讨论中的提议也不因写入本书而自动成为定理。}

读者可以先读序言与“思想前史”，看见这套体系从“路径原创”、语言适配与 Return 的问题中怎样生长，再读宪法与两次判例。条款随后以“来源 + 压力测试状态”挂回这些判例，AI 执行模板则从已经取得资格的规则派生。全文的完整性属于理由档案，实际协作的简洁属于执行层。两者不再各自漂移。

\vfill
\noindent{\small\color{muted}版本约定：宪法规定最高优先次序；判例保存规则形成与修订的理由；条款必须标明来源与压力测试状态；AI 模板从这些条款派生。第三、四条宪法在 v0.3 中暂定，等待结构不同的后续判例复验。任何一层都没有免于检验的特权。}
\tableofcontents
\mainmatter
\standalonechapter{序言：理解的责任与创造的自由}
我们常常以为，理解一个理论，就是能够说出它的定义，复述它的结论，跟随一遍证明。然而，有一种不满足仍会留下：即使每一句都听懂了，整个理论为什么偏偏长成这个样子，似乎还没有真正进入自己的认识。

这种不满足不是知识不足的另一种名称。有时，一个人已经读过许多文献，能够熟练使用术语，却仍然只能沿着别人已经铺好的路行走。遇到熟悉的对象，他知道应当调用哪一个定理。遇到没有现成名字的问题，他却不知道该从哪里开始。知识积累与认知自由之间，并没有一条自动通行的道路。

闭依开生法正是从这条道路的缺口出发。它要同时完成两个任务：向后，恢复已有知识的成立方式；向前，使这些成立方式不再成为新知识的唯一模板。前一个任务叫“闭依”，即 Dependency Closure。后一个任务叫“开生”，即 Generative Independence。

\lead{理解，是对已有依赖负责；创造，是不让已有依赖占有未来。}

这两句话之间有一种必要的张力。理解要求我们忠实于已经发生的推理，不把别人付出的证明工作抹成一个轻巧的结论。创造却要求我们不把已经发生的历史误认成未来必须服从的法则。一个定理通过某种几何实现出来，并不意味着那种几何就是定义本身。一个证明使用了某种递归，并不意味着同一个结论只能由递归产生。

因此，闭依不是把过去背得更熟，开生也不是对过去摆出独立姿态。前者要问“究竟什么使它成立”，后者要问“这些作用还能由什么承担”。只有前者足够清楚，后者才不至于沦为换名。只有后者始终开放，前者才不至于变成对历史形式的忠诚。

作者的认识习惯为这套方法提供了生成的起点：“先识全貌，再求甚解；先建立形式上的共鸣，再理解具体的证明。”这里的全貌不是一段概述，而是对对象、关系与作用方向的整体感受。形式上的共鸣也不只是两个名词相似，而是尚未逐项验证之前，已经隐约听见它们在某种关系安排上互相呼应。

这很像理解一首乐曲。听者先感受到整首作品的呼吸，再辨认主题怎样出现、怎样变形、怎样在不同声部中获得新的意义。技术分析不会使第一次聆听成为多余，但第一次聆听也不能替代对困难段落的核验。直觉让一个整体成为可能，细读决定这个整体究竟站不站得住。

DCGI 不试图消灭这样的直觉。它为直觉保留提出结构的权利，也为反例保留修改结构的权利。它反对的，一面是“我感到相容，所以已经相容”，另一面是“还没有完整语言，所以不能开始思考”。认识可以先于命名发生，结论却必须晚于它所需要的检验。

这套方法并不是一次设计完成的。第一次压力测试使它看到：一个概念不能被它最熟悉的实现占有。第二次压力测试则暴露出另一种危险：当方法论开始自我运转，它也会占有对象。审计者可以知道许多一般理论，却忘记眼前的问题已经给出了更具体、更直接的结构。于是，增加分析不再意味着增加理解，而可能只是增加偏离。

作者最后的纠正很短：“不要生搬方法论，大道至简。”这不是替不严谨开脱，而是要求方法回到其职责。简洁不是省略必要依赖，严谨也不是把一切可能的困难都带进来。\textbf{需要留下的，是当前结论真正依赖的东西；需要删去的，是方法自己制造的噪声。}

由此，DCGI 才有了今天的层级。宪法规定最高优先次序，判例保存规则为何形成及怎样被修正，条款据此取得各自的来源与压力测试资格，AI 模板只负责把已经取得资格的约束压缩成可执行协作。整套体系最终要走向一种朴素的成熟：面对知识，能够承担理解的责任；面对未知，保有重新生成的自由；面对一个本来清楚的结构，知道何时不再说方法论。


\standalonechapter{思想前史：路径、语言与返回}
\label{ch:prehistory}

DCGI 的正式名称是在数学研究的压力下形成的，但它并不是从一张空白纸上开始。更早的一组文章已经反复触及同一个认识论难题：一个人怎样在大量现成语言之间保持自己的问题，又怎样在不拒绝外部知识的前提下，使理解真正成为自己的历史？这些文章没有预先给出“Dependency Closure”或“Generative Independence”这样的术语，却逐步准备了它们后来能够成立的地面。

这段前史值得保留，不只是为了交代思想来源。DCGI 主张理解必须对依赖负责，那么它也必须对自己的生成依赖负责。若一个方法论把自己的形成史删成四条整齐原则，它会重演自己反对的事情：让最终概念遮蔽问题怎样一步一步长出来。

同时，这一章也必须守住边界。以下文章属于认识论与跨学科反思，它们提供问题、区分和语言，不替代数学证明。Just infiniteness 与 orbit equivalence 的两次压力测试后来所做的，正是把这些较为开放的直觉送进一个更严格的环境，看它们究竟能保留多少、必须修正什么。

\section{概念最危险的时候，是它来得太早}

《问题先于概念：AI时代为何要保护自己的理解路径？》从一个极简单、也极容易被忽略的经验开始：一个问题刚刚发生时，我们常常还不知道自己究竟看见了什么。此时，一个成熟概念若过早出现，会带来强烈的舒适感——现象终于有名字了，文献也已经在那里——但这种舒适可能只是把问题提前关闭。\citep{peng-problem}

概念当然不是敌人。没有概念，复杂经验无法被压缩、比较和累积。危险只在于顺序：如果命名先于问题自己的分化，那么研究者很容易开始只看见该概念允许他看见的东西。现实于是从一个尚未解决的问题，变成某个理论的案例。

这里第一次出现了后来 DCGI 的一个核心倾向：\textbf{能够暂时不知道“这叫什么”，是一种积极的认识能力。}它要求我们在命名之前多停留一会儿，让对象先暴露自己的差异，让问题先拥有形状。所谓“延迟命名”，并不是反理论，而是为理论选择保留一个前置阶段。

这也解释了为什么 AI 时代尤其需要这项纪律。搜索、生成和跨语言检索越强，现成概念到达问题的速度越快。过去，一个研究者可能因为不知道某个术语而被迫绕路；今天，他可能在问题尚未成形时就被十种相关 framework 包围。知识获取的摩擦下降了，\emph{被现成语言提前吸收的风险反而上升}。

后来 DCGI 所说的“问题先于语言”“机制先于术语”，都可以在这里看到雏形。但当时更重要的还不是规则，而是一种感受：真正值得保护的，不是“我自己的说法”，而是\textbf{问题在尚未被归类时仍然能够继续说话的时间}。

\section{理解有路，而路不能预先占有理解}

《理解有路，真理无路》把“路径原创”推进了一步。文章借克里希那穆提“真理是无路之国”的命题追问：如果人类理解必然从一个有身体、有语言、有历史的位置开始，那么理解不可能是从无位置处一次性获得对象本身。理解必然有路径。\citep{peng-path}

但“理解必然有路径”并不推出“存在一条普遍正确的路径”。恰恰相反，如果路径在理解之前已经被完全给定，那么研究者真正做的就不再是理解，而是按照既定终点塑造自己的经验。文章因此把三层意思连在一起：理解必然有路径；真正的路径只能在理解过程中生成；因此，没有一条能够先于理解而存在、并替所有问题规定终点的普遍道路。

这一步与后来 Generative Independence 的关系很深。已有理论当然可以成为道路的一部分，甚至可以极大缩短新的探索。但一条历史上成功的道路不能因为成功，就获得对未来问题的所有权。真正需要保存的是：\textbf{路径可以继承，路径权不能继承。}

从这里看，“原创”也获得了不同于“没有前人说过”的含义。路径原创不是故意绕开文献，更不是把已有概念换成自己的词。它首先意味着：当前问题怎样组织已有知识，必须由当前问题本身决定。新的路径完全可以经过旧桥，却不能因为桥已经存在，就假定所有人都必须去同一个目的地。

这后来会在 just infiniteness 的压力测试中得到严格版本：定义、机制与历史实现必须分开。一个概念第一次在 rooted-tree machinery 中变得清楚，并不意味着树就是概念本身。所谓“不让已有依赖占有未来”，并不是反历史，而是拒绝把\emph{曾经如此}偷换成\emph{只能如此}。

\section{语言不是阶梯，而是与问题相配的坐标}

《语言之间不是等级关系，而是问题—语言之间存在适配关系》进一步修正了一种非常常见的认识冲动：既然世界只有一个，是否终有一种描述语言比其他语言更基础、更真实，最后能够把其余语言全部吸收？文章给出的回答不是简单的多元主义，而是一个更操作性的判断：任何描述都通过舍弃一些差异来突出另一些差异，因此真正应当评价的，不是哪种语言“最高”，而是哪种语言保存了当前问题所需要的结构。\citep{peng-language}

这条原则有两个方向。第一，一个理论失败，有时不是因为它错误，而是因为我们让它回答了它并不适合回答的问题。第二，一个理论在自己的适用域中极其成功，也不能因此无限扩张解释权。语言越成功，人越容易忘记自己正在使用一种语言，把地图当成土地，把模型当成对象。

这正是后来“对象主权”的思想前身。对象主权并不是拒绝抽象语言，而是拒绝语言凭借自身成熟程度取得统治地位。数学、动力系统、群论、范畴论、心理学与哲学各自能够保留不同结构；它们之间的桥梁需要被证明或至少被明确限定，而不能靠术语迁徙建立。

从 DCGI 的角度看，所谓 projection before absorption 也因此不是一个信息整理技巧。它是在问：这个外部理论进入我的问题时，究竟照亮了哪里，又遮掉了哪里？如果它只回答相邻问题，就应当让它停在相邻位置，而不是因为它拥有漂亮的语言就重写原问题。

\section{Projection 使路径出现，Return 保存来路}

《为什么我们必须绕远路？》把“路径”从认识论进一步推进到结构生成。文章从误解与投影出发指出：投影当然意味着信息损失，但有限视野并不只有消极意义。正因为一个位置不能一次性占有整体，理解才需要沿着路径展开。\citep{peng-detour}

这里形成了一条重要链条：
\[
\text{Generation}
\longrightarrow
\text{Differentiation}
\longrightarrow
\text{Projection}
\longrightarrow
\text{Path}
\longrightarrow
\text{Return}.
\]
生成使结构出现，分化使同一发生史进入不同位置，投影让每个位置只有有限视野，路径因有限性而形成，而 Return 则使已经分化的结构重新发现来路。

但文章最重要的限定，是\textbf{Return 不是 Rewind}。若一个系统经过变化之后只是恢复到未经变化的原状，那么中间的历史实际上被抹掉了。真正有意义的 return 更接近
\[
O\longrightarrow \text{Differentiation}
\longrightarrow \text{Projection}
\longrightarrow \text{Path}
\longrightarrow O',
\]
其中 $O'\neq O$，却仍与 $O$ 保存某种可以说明的结构连续性。返回不能取消差异，而必须把路径一起带回来。

这一点后来成为 Dependency Closure 的深层背景。闭依不是倒退到“在我尚未学习这些理论之前，问题原本是什么样”，也不是恢复一个未经文献影响的纯净自我。那种纯净状态既不可得，也没有必要。闭依要做的是：在已经发生过学习、误解、修订和吸收之后，仍然能够说明\textbf{我怎样来到这里}。理解成为自己的历史，不是因为它没有受到影响，而是因为这些影响可以被重新辨认和组织。

因此，provenance 在 DCGI 中不是档案学的附加要求。它对应一种更深的 returnability：一个认识过程若完全无法恢复自己的来路，就很难区分哪些结构已经成为自己的理解，哪些仍只是被临时借来的 scaffold。

\section{共同来源并不保证可返回，兼容也不等于同源}

接下来的《Return 的边界在哪里？》对前一篇文章主动加压。若 Monogenesis——共同生成——能够帮助组织 return，那么是否只要两个结构曾经同源，就必然还能返回？文章明确否定了这一步：共同来源是历史命题，可返回性却是当前结构命题。两个结构可以确实来自同一源头，却在长期分化后失去任何可恢复路径。\citep{peng-return-boundary}

于是，Similarity、Monogenesis 与 Compatibility 必须拆开：Similarity 比较当前结果有多像，Monogenesis 追踪是否拥有共同生成史，Compatibility 则问已经不同的结构能否在某个约束或接口下进入共同结构。它们可以重合，也可以彼此独立。

这一区分与 DCGI 后来的 realization discipline 直接呼应。形式相似不能自动升级成共同机制，共同来源不能自动升级成当前可实现的连接，而局部 compatibility 也不能自动升级成 global realization。真正的桥梁需要额外结构。

文章借 pushout 的形式启发进一步提出：即便旧的共同来源已经不存在，新的共同结构仍可能通过一个后来形成的 interface 被创造出来。这里重要的不是把哲学中的 Return 等同于范畴论 pushout，而是获得一个新的问题：\textbf{整合是否一定依赖回到同一中心，还是也可以依赖后来构造出的共同可组合条件？}

这为 DCGI 的“开生”准备了另一面。Generative Independence 并不总是寻找更古老的本质；它有时是在已有分化之后寻找新的可组合结构。新的机制可以忠实于旧问题，却不必复制旧历史。

\section{中心消失以后，整合不能跟着消失}

《当中心消失之后，结构如何重新整合？》继续追问：如果三阶意识要求观察者本身进入观察，一个固定主体不再拥有最后解释权，那么去中心化会不会最终只剩碎片？\citep{peng-center}

文章的回答不是恢复一个更高的“超级主体”，而是把统一从“共同中心”转向“共同组织条件”。一个 interface 可以使不同位置保持差异，同时获得组合、转换与继续的可能。统一于是可以是关系性的，而不是封顶性的。

这个思想后来在 DCGI 中表现为两项互相制约的要求。一方面，任何外部 framework 都不能取得 final adjudicative authority；问题本身始终保留修改框架的权力。另一方面，这并不意味着各种解释只是平行存在、互不相干。Dependency Closure 要求它们进入一个可以追踪的依赖结构，使“没有最终中心”不等于“没有组织”。

所以，DCGI 所追求的不是一种认知上的无政府状态。它反对的是一个位置永久封顶，而不是反对结构。事实上，越没有最终权威，越需要明确 dependencies、provenance、interfaces 与 failure boundaries，否则“开放”很容易变成无法检验的任意性。

\section{AI 可以重新接通语言，却不能接管问题}

《AI可否克服语言分裂，重建人类巴别塔？》把上述问题推进到 AI。AI 可以大幅降低语言之间的转换成本：翻译论文、解释术语、改写表达、连接跨学科材料。技术上，巴别塔式的“听不懂对方在说什么”正在迅速减弱。\citep{peng-babel}

但语言可互译并不意味着观察位置已经消失。如果 AI 把所有 vocabulary 都接通，真正的问题反而更加裸露：不同的人与不同学科是否愿意承认彼此位置中的不可还原部分？一种统一语言是否又会以技术便利的名义成为新的共同中心？

DCGI 因此不给 AI problem sovereignty。AI 可以做非常强的检索、比较、重述、压力测试与反例搜索，却不能因为它知道更多语言，就默认知道问题应该被哪种语言组织。它应当帮助研究者提高必要的分辨率，而不是以“更 general”的名义把问题迁往自己更熟悉的环境。

这条约束在第二次压力测试中得到了一次几乎教科书式的失败：一般 orbit equivalence 的信息损失被机械带入一个更具体的 AF-by-discrete 语境，方法知道得更多，反而看得更远离对象。正因为这次失败真实发生过，“对象主权”与“大道至简”才从漂亮原则变成了暂定宪法条款。它们目前仍等待结构不同的后续判例独立复验。

\section{从思想前史到 DCGI：继承的是问题，不是术语}
回看这些文章，DCGI 并没有把 Monogenesis、Return、Interface、Projection 或“路径原创”直接写进宪法。它继承的是更深的问题：怎样让问题在命名之前保持开放，怎样让语言服务于问题而不垄断对象，怎样保存已经走过的路径，又怎样在取消最终中心以后仍然维持可检验的组织。

进入数学研究以后，这些问题被迫获得更硬的资格。Just infiniteness 要求区分 definition、mechanism 与 realization；orbit equivalence 则迫使讨论辨认自然语境，并让方法论承认：增加分析层级并不总是增加理解。思想前史提供问题敏感性，压力测试决定这些敏感性究竟能保留到哪里。

因此，思想前史在理由档案中只承担生成史角色。它解释为什么后来会有某些规则，却不替这些规则取得数学证明或永久宪法资格。DCGI 对自己的前史也实行同样的纪律：保留来路，但不让来路预先占有结论。

\chapter{宪法}
\label{ch:constitution}
宪法只规定最高优先次序。前两条是 DCGI 的定义性轴线：闭依与开生。第三、四条来自第二次压力测试，目前列为\textbf{暂定宪法}。暂定不是削弱其实际约束力，而是保存其证据资格：只有当结构不同的后续判例再次独立要求同一原则时，才考虑将其转为确认宪法。同类问题的重复成功不能用次数代替结构独立性。

\Article{闭依原则}
\lead{理解，是对已有依赖负责。}
对一个已有知识对象，不能只继承结论，而应在当前问题所需的范围内，恢复它的对象、定义、成立机制、适用条件与失效边界。所谓闭依，不是把每一个基础定理重新证明一遍，而是不让关键依赖在“我大概知道”中失去踪迹。

我们至少应当分清：一个条件写在定义中，是因为它规定了所讨论的对象，还是因为它便于某种证明？一个局部构造能够推广到整体，是依靠紧致性、对称性、有限生成，还是一个另外证明的延拓结果？一个结论改变表示之后仍然成立，是已有不变性定理，还是仅仅因为记号没有改变而被默认了？

闭依的完成标准是\textbf{依赖关系透明}。重要节点都有明确身份：已经证明的、作为前提接受的、由文献调用的，以及尚未闭合的。读者应能沿着这些节点返回结论的来源，而不是在一连串熟悉词语之间寻找通路。

这种透明也包含对历史的理解。我们需要知道一条道路怎样开辟，却不能把历史先后当作逻辑必然。理解“为什么最早如此实现”与理解“为什么必须如此实现”，是两个不同的问题。闭依必须将它们分别回答。

\Article{开生原则}
\lead{创造，是不让已有依赖占有未来。}
对已有知识负责，并不意味着未来必须复制它的成立方式。一个理论可以告诉我们某种机制足以产生某个性质，却未必证明它是唯一的机制。开生要求在理解之后重新释放这个可能性。

这里有三件必须保持区别的事：定义规定什么算作目标对象，机制解释目标性质怎样发生，实现给出机制具体工作的场所。它们可能在一个经典例子中紧密重合，但认识上的重合不构成逻辑上的同一。

一个来自有根树的证明，可能依赖的不是“树”这个名称，而是局部支撑、共轭传播与有限指数控制。离开树时，需要重新寻找这些作用能否被承担，而不是为每个树上名词发明一个新语境中的同义词。反过来，新的问题也可能需要完全不同的机制。开生并不要求忠实保留旧证明的全部结构，只要求新证明承担自己声称的结论。

开生的完成标准是\textbf{已有知识成为材料，而不再成为模板}。它不是必须推广的命令，也不是必须原创的表演。发现一种新实现是结果，找到一个不可跨越的结构障碍也是结果。还没有找到时，准确地留下问题，比宣布一个空泛的新框架更接近创造。

\Article{对象主权（暂定宪法）}
\lead{方法论服从对象，而不是对象服从方法论。}
每个问题都带着自己的语境。讨论开始时已经固定了哪些结构，允许改变哪些表示，正在比较什么对象，这些事实不能在审计中被悄悄擦掉。

因此，调用一般理论之前，应先看清当前的环境、对象、允许的变换，以及随问题一起给出的背景结构。本文称这种工作为“确认自然语境”。它不是给研究设置围墙，而是保证每一次离开原语境都是主动的选择，而不是未被察觉的偏移。

若一个问题本来讨论带有共同 AF 核的作用，却在论证中变成任意 Cantor 作用的轨道等价，那么审计的对象已经改变。一般语境中的困难未必错误，但不能直接用它否定原问题中的结构。本文将这种偏移称为“语境漂移”，英文为 \emph{ambient-category drift}。

对象主权同样约束作者与 AI。熟悉某种语言不意味着有权让所有问题进入那种语言。作者的直觉也不能因此取得事实豁免权。\textbf{优先的是对象所具有的结构，不是任何一方解释结构的权威。}

对象主权还要求一个对称的判别标准：\textbf{质疑必须指出它作用在原论证的哪一步；驳回质疑也必须指出哪一个结构理由使那一步成立。} 若审计者说“这里可能有 gap”，却无法给出受影响的 implication、definition 或 estimate，这个疑点暂不进入正式审计。若作者只说“这里没有问题”，却没有指出缺失的结构桥梁，同样不能据此关闭质疑。作者身份与 AI 怀疑都不是证明资格。

开生当然允许改变语境。它与对象主权并不矛盾：前者要求有意识地探索另一种实现，后者禁止把这种改变伪装成原问题本来如此。
\claimstatus{暂定宪法。直接来源为判例二；本条下“质疑与驳回均须定位具体步骤”的对称判别标准同样为暂定；等待至少一个结构不同的后续判例独立复验。}

\Article{至简原则（暂定宪法）}
\lead{方法论只增加必要的分辨率，不为自身制造复杂性。}
一个方法的价值不由条款数量、术语密度或怀疑强度决定，而由它是否使当前问题更清楚决定。能直接识别的结构，应当直接使用。只有当某个歧义、依赖缺口或实现问题确实阻碍理解时，才需要进一步拆分。

至简不是尽量少写字。一份面向读者的文档可以很详尽，因为它需要说明背景、展示转折、保存可核查的细节。至简所要求的，是\textbf{推理不承担与结论无关的负担}。解释可以充分，逻辑不必绕远。

它也不是把“显然”当作免检通行证。真正的简洁能够指出必要结构为何足够。若一条被省略的依赖实际上决定结论，补出它是在减少误解，而不是制造复杂度。反之，若一个新增层级没有改变任何判断，也没有关闭任何缺口，就应当考虑把它删掉。

当方法开始遮住对象，至简原则要求停止那一套操作，回到最短有效的解释。它否决的是多余程序，不是真实性与证明义务。所谓“大道至简”，在这里不是一句审美口号，而是一条工作约束：\textbf{让复杂性来自对象，而不是来自对方法的忠诚。}
\claimstatus{暂定宪法。直接来源为判例二；等待至少一个结构不同的后续判例独立复验。}

\chapter{判例与压力测试}
\label{ch:cases}
一套方法可以在说明中显得无懈可击，却在真实对象面前暴露盲点。因此，本编不只记录方法怎样帮助研究，也记录它怎样误导研究。两次测试的价值并不对称：第一次打开了概念与实现之间的距离，第二次则限制了方法继续制造距离的冲动。

下文中的数学只保留理解转折所需的内容。经典结果给出出处，讨论中的候选机制明确称为候选。对话中的短句按原意保留，其余内容按逻辑重述。

\section{判例一：一个概念能否离开它最熟悉的实现？}
\label{sec:ji}
\subsection*{起点：从 just infiniteness 追问成立方式}
第一次测试并不是要求解释一个陌生术语。真正的问题是：just infiniteness 为什么值得研究？它在有根树上怎样被证明？离开这套熟悉的几何以后，是否还能找到由 tile 或群胚结构直接产生的机制？

这三个问题分别指向概念、证明与构造。如果不把它们分开，所谓理解往往会变成：知道一个经典群具有这个性质，知道它作用在树上，于是把树当作该性质本身的组成部分。

先看定义。一个群 $G$ 称为 just infinite，如果它是无限群，而且每个非平凡正规子群都有有限指数，即
\[
  \{1\}\ne N\triangleleft G\quad\Longrightarrow\quad [G:N]<\infty.
\]
等价地，每个由非平凡正规子群得到的商群都是有限群。这一定义只涉及群与正规商，并没有指定有根树。相关结构理论及其分支群背景可见 \cite[Introduction, Chapter~5]{bgs03}。

最简单的检验甚至不需要复杂例子。无限循环群 $\mathbb Z$ 的非零子群都是 $n\mathbb Z$，相应指数为 $n$，因此它满足定义。$\mathbb Z^2$ 则不满足，因为对 $\mathbb Z\times\{0\}$ 取商仍得到无限循环群。一个无限单群也满足定义，因为它没有非平凡的真正规子群。这些直接算例已经足以说明：\textbf{同一个商意义上的性质，不必由同一种几何机制产生。}

关于历史，讨论从 Grigorchuk 的工作进入这个问题，但不应因此把概念史简化成“始于树上”。Wilson 在 1971 年已研究每个真商都有限的群，后来的分支理论对这类对象的结构与实现作出了重要推进。这里追问的是这条研究路线解释了什么，而不是用一个熟悉名字代替完整的历史。\cite[Introduction]{bgs03}

\subsection*{第一处转折：把树上的证明拆回它所做的事}
在有根树 $T$ 上，一个顶点 $v$ 决定其下方的子树 $T_v$。刚性稳定子 $\rist_G(v)$ 由那些在 $T_v$ 之外逐点恒等的群元素组成。它的意义不只在于“稳定一个顶点”，而在于把非平凡作用限制在一个可分离的局部区域内。

对同一层的顶点，可以将这些局部子群放在一起。由于不同子树互不相交，它们的作用能够交换。随后，群在层上的传递性使一个区域中的信息传播到其他区域。分支群的经典正规子群论证正是沿着这种局部化与传播工作：非平凡正规子群包含某一层相应刚性稳定子的导出子群，而有限指数结论还需要相应的有限性条件。\cite[Theorem~5.2, Lemma~5.3]{bgs03}

这条思路可以压缩为
\[
\text{发现非平凡作用}
\longrightarrow
\text{局部化}
\longrightarrow
\text{共轭传播}
\longrightarrow
\text{有限指数结论}.
\]
最后一步不能被“传播”这个词自动包办。获得许多局部作用，仍须说明这些作用合起来为什么足够大。正是在这里，定义、机制与实现开始分离：树提供了一种局部区域的组织方式，局部化与传播是证明动作，just infiniteness 才是目标性质。

\subsection*{一个可以直接核查的局部计算}
为了看清哪些内容真正依赖树，可以把局部化写在更一般的空间上。设 $\Gamma\leq\operatorname{Homeo}(X)$，采用右作用，并令
\[
\Rist_\Gamma(U)=\{g\in\Gamma:xg=x\text{ 对所有 }x\notin U\}.
\]
这是只在 $U$ 内允许非平凡作用的子群。设 $N\triangleleft\Gamma$，且 $n\in N$ 满足 $U\cap Un=\varnothing$。取 $a,b\in\Rist_\Gamma(U)$，约定 $[x,y]=x^{-1}y^{-1}xy$。由于 $n^{-1}a^{-1}n$ 支撑在 $Un$，它与 $a,b$ 交换，因而
\[
 [[n,a],b]
 =[(n^{-1}a^{-1}n)a,b]
 =[a,b].
\]
左边属于 $N$，所以 $\Rist_\Gamma(U)'\leq N$。

这个计算没有使用树。它使用的是不相交的支撑与正规性。它解释了为什么离开树后仍可以寻找局部化机制，却没有说明怎样从一个 $U$ 扩散到足以控制整个群的子群。这后一件事，仍要由新对象自己的结构承担。

\textbf{第一次测试的进展，不是宣布“树可以去掉”，而是知道哪一步已经不需要树，哪一步仍在等待新的实现。}

\subsection*{第二处转折：不能只给一张抽象格}
讨论随后遇到一个更直接的要求：仅仅说“研究正规子群格”还不够。定义本来就是关于正规子群的。若想发现新的证明机制，必须看见具体子群怎样产生、怎样嵌套、怎样交换、怎样在共轭下移动。

因此，要求转向了显式的子群或局部结构格。在树上，顶点与子树已经提供了一套可见的区域系统。在 Cantor 动力系统中，一个自然起点是 clopen 集，即既开又闭的子集，以及上面定义的 $\Rist_\Gamma(U)$。若 $U\subseteq V$，就有 $\Rist_\Gamma(U)\leq\Rist_\Gamma(V)$。若 $U\cap V=\varnothing$，两个子群逐元素交换。这些都是可以直接检验的关系。

但将 $U$ 配给一个支撑子群，尚未得到足以证明 just infiniteness 的结构格。还需要知道这些子群是否非平凡，是否在给定群中实际存在，以及它们的传播能否给出有限指数。若所需局部变换只存在于更大的拓扑全群中，而不在正在研究的子群中，就不能将前者的能力记在后者名下。

这条区分对 AF-by-discrete 研究尤为重要。一个群胚含有 AF 核，不等于某个给定作用群已经含有 AF 核的整个全群。\textbf{环境拥有的局部对称性，与所选群实际能够实现的局部对称性，必须分别核对。} 这使“显式结构”成为实质要求，而不只是把抽象格画得更形象。

\subsection*{第三处转折：承认不止一种机制}
讨论一度容易把全部希望放在“寻找树外的分支机制”上。但 just infiniteness 的目标性质本来比这条路线更宽。已有结构理论将分支群与由遗传 just-infinite 群或无限单群所构成的其他情形联系起来，相关结论包含有限指数正规子群及其直积分解等条件，不是简单地把所有群贴成三个标签。\cite[Theorem~5.6]{bgs03}

在这次测试中，它的作用是一张机制阅读地图：有的刚性由局部支撑的复制与传播产生，有的由有限指数子群中持续保留的商刚性产生，有的则与单纯性有关。认识这张地图以后，就不能再把“找到树外分支”当作实现目标的唯一方式。

在 Cantor 动力系统的研究中，无限单群与全群结构本来就提供了另一条具体路线。例如 Nekrashevych 的工作构造了中间增长的有限生成无限单群。\cite{nek18} 但这里也不能把“有限指数单纯核心”四个字当作自动结论。一个很短的检验就能看出需要补什么。

设 $S\triangleleft G$ 是无限非阿贝尔单群，$[G:S]<\infty$，并且 $C_G(S)=\{1\}$。对任意非平凡正规子群 $N\triangleleft G$，$N\cap S$ 要么等于 $S$，要么平凡。前者直接给出 $[G:N]<\infty$。后者则有 $[N,S]\leq N\cap S=\{1\}$，从而 $N\leq C_G(S)=\{1\}$，矛盾。因此这些条件确实推出 $G$ just infinite。

若删去中心化子的控制，结论就不能照搬：$S\times C_2$ 含有有限指数的无限单群 $S\times\{1\}$，却也含有指数无限的非平凡有限正规子群 $\{1\}\times C_2$。这个近失例的意义很具体：它不是否定单纯核心路线，而是指出这条路线最后必须承担的依赖。

\subsection*{AF-by-discrete 中的候选方向与停止线}
回到作者的对象，tile 的有限层次、clopen 区域、局部交错群与共轭传播，都可能为新的实现提供材料。近全群研究中涉及的局部实现与同步限制，也提示了哪些信息需要被显式跟踪。

这些是候选方向，不是一条已经完成的“非树分支定理”。一次真正的构造仍须展示具体子群族，证明所需局部作用属于目标群，解释传播为何有效，并给出有限指数或相应刚性的收口。若只把刚性稳定子换成“tile 稳定子”，却没有完成这些工作，就只是改变了语言。

第一次测试因此既打开了道路，也留下了停止线。它获得了更准确的问题，却没有强迫问题立刻变成一条总定理。\textbf{开生不是必须推广，而是在能够重新选择机制之后，仍然诚实地区分已实现与待实现。}

\subsection*{本案确立的原则}
本案最重要的结论是：定义不由第一次成功的实现独占。要使一个概念走出熟悉环境，先恢复证明中各个对象承担的作用，再在新对象中寻找显式实现。抽象性质与分类结果提供方向，具体子群与局部计算才提供抓手。

由此进入 DCGI 的，是“定义、机制、实现分离”“依赖敏感的理解深度”“实现债务”与“开生不等于强行推广”。这次测试证明的不是一种新的普适群论，而是一种更可靠的提问能力。

\subsection*{判例登记：Just infiniteness}
\begin{tabularx}{\textwidth}{@{}p{30mm}X@{}}
\toprule
触发事件 & 追问 just infiniteness 为何在 rooted-tree 语境中出现，以及离开树后能否重建其证明机制。\\
原判断 & 容易把 rigid stabilizer、branching 与 just infiniteness 本身捆在一起。\\
偏离方向 & 将经典实现误当成 intrinsic definition，或只给抽象正规子群格而没有显式局部结构。\\
纠正依据 & 回到 just infinite 的商刚性定义，并逐步抽取局部化、传播与有限指数控制的实际作用；双交换子计算给出一个可直接核验的局部机制。\\
影响条款 & 定义--机制--实现分离；认知深度服从依赖；实现债务；开生不等于强行推广。\\
当前状态 & 定义与引用的结构结果属于\textbf{文献调用}；文中双交换子小计算属于\textbf{已证}；AF-by-discrete 中新的 just-infiniteness 机制仍是\textbf{候选/待闭合}。\\
\bottomrule
\end{tabularx}


\clearpage
\section{判例二：当审计离开了它应当理解的对象}
\label{sec:oe}
\subsection*{起点：一个具体的传递问题}
第二次测试开始于一个并不宽泛的问题。给定由 tile inflation 描述的有界类型逆半群 $G_0$，考虑其拓扑全群中的有限生成子群 $G\leq\FF(G_0)$，并在相应的轨道比较语境中研究：$G_0$ 的有限不可压缩性怎样传递到 $G$？相关论文的引言明确提出了这一传递，而其主要增长结果另有展开性、线性重复性等假设。\cite[Introduction, Section~4]{kua25}

为了理解问题，先保留四个对象。$\BB$ 是给定的 Bratteli 图，$\Omega(\BB)$ 是其无限路径空间。两条路径共尾，是指它们只在有限多个初始位置不同。由此得到尾等价关系，其 AF 群胚记为 $\TT$。tile 提供有限近似，持续边界点则记录那些在不断细化中仍未被内部化的边界位置。

“回返”在这里不是抽象的“又发生一次”。它有指定的路径、首尾边界标签与内部运行条件。论文对 return word 与 incompressibility 的定义正是以这套 tile 与边界约定为基础。\cite[Definitions~2.17--2.18]{kua25} 因此，讨论是否传递时，首先应追踪所用的几何结构，而不是仅仅观察两个定义中都出现了同一个英文词。

\subsection*{第一处偏移：把一般事实放在了错误的位置}
审计最初调用了一条一般性的警觉：普通 orbit equivalence 记录轨道如何对应，不自动记录全部局部胚、箭头拓扑或生成元标记。对两个作用而言，知道哪些点属于同一轨道，与知道某个变换在邻域中怎样实现，本来就是不同层次。

作为一般提醒，这一区分有意义。问题在于，它很快取代了眼前已经给出的完整语境。讨论不只有一个裸的轨道对应，还包含 $G\leq\FF(G_0)$，以及正在共同追踪的 AF 尾结构。审计者没有先看这些信息到底做了什么，就把问题改写成“普通 OE 是否保存精细组合性质”。

于是，一条具体的传递问题被扩展成了许多新的议题：胚是否被保存，词语言是否被保存，回返怎样穿过编码边界，是否需要独立的局部重编码定理。这些问题有的在别处完全合理，但是否是当前问题的承重点，尚未被建立。

这就是后来被命名为语境漂移的错误。它不一定表现为某句一般知识说错了，而更常表现为：\textbf{一般知识没有回答当前问题，却因其正确与熟悉而占据了当前问题的位置。}

\subsection*{一个原本就存在的局部事实}
$G\leq\FF(G_0)$ 的含义，是每个 $G$ 元素在每一点附近都与某个 $G_0$ 元素一致。\cite[Section~4, p.~14]{kua25} 对固定的有限生成集 $A\subseteq G$，紧致性允许为每个 $a\in A$ 选择有限 clopen 分割，使
\[
  a|_{U_{a,i}}=F_{a,i}|_{U_{a,i}},\qquad F_{a,i}\in G_0.
\]
固定 $G_0$ 的有限标签集后，这些有限多个局部代表的字长有共同上界。

这是一个直接来自定义与紧致性的事实。它说明本题并不是只知道逐点的轨道可达性。另一方面，这个事实也不应被夸大成全部组合传递的替代证明：常数属于固定的生成集与这些局部代表，并非无条件地控制所有可能表示。

真正恰当的使用，是把这条局部事实放回对象的已有结构中。它既不应被一般 OE 的警觉抹掉，也不必被命名成一套越来越庞大的新机制。

\subsection*{第二处偏移：把没有证明说成可能不成立}
审计随后用了过强的措辞：$G_0$ 有限 IC “不一定推出” $G$ 有限 IC。作者追问了一个很具体的问题：有没有反例？

这一问把讨论从论证姿态带回了数学资格。若要否定一个蕴含，必须给出满足全部假设、却不满足结论的对象。仅仅解释“为什么我眼下没有把证明接起来”，还不能承担这个否定。

讨论中曾尝试用二面体作用与碎片化作用提供候选：若某个碎片化群与原作用具有相同轨道，再加入原作用中的无限阶元素，是否就得到反例？但这条思路没有核实新增元素属于碎片化作用的拓扑全群。缺少这一项，它就没有满足原问题的全部假设，因而没有取得反例资格。这里不需要先证明这个候选永远不能成功，才能承认当时尚未构成反例。

审计者因此撤回了过强判断。应该留下的是“尚未在所检查的论证中建立这一传递”，而不是“已经证明传递失败”。这是结论强度匹配在否定性判断上的一次实际检验。

然而，作者进一步要求证明时，检索又反复返回同一声明、同一定义与相近的副本。检索量增长了，论证却没有相应推进。它暴露出一个更日常的错误：\textbf{当任务已经要求推理，继续收集同一种背景材料并不等于继续完成任务。}

\subsection*{术语的小岔路：reduced 并不是熟悉的那个 reduced}
这段审计中，作者还追问：“这个 reduced word 你有没有明白是什么意思？”这个问题之所以重要，是因为一个熟悉的词也会造成语境漂移。

在所讨论的定义中，$w=s_1\cdots s_m$ 被称为 reduced，需要非恒等字母及其相邻配对满足指定的边界条件：相邻字母对应不同无限边界点上的边界标签，或者至少有一个是 finitary。它并不是通常的自由约化，也不自动表示群元素的最短词。\cite[Definition~4.2, Remark~4.5]{kua24}

其作用是区分证明面对的不同边界行为。原证明分别处理 reduced 情形与含非 reduced 子词的情形，后一情形通过适当层次的 section 与诱导穿行组织论证。读者若把“reduced”按熟悉的代数意义代入，就会追问一个并非该定义要回答的问题。

Grigorchuk 例子提示了一个方法论问题：当预期使用的 reduced 情形不能覆盖所需轨迹时，应回到定义与证明的实际用途，检查另一情形如何处理。这个例子说明核验定义用途的必要性，本身不构成新的普适数学断言。

这条岔路进一步说明：对象主权不只约束宏观的环境，也约束一个技术词的具体定义。

\subsection*{决定性的纠正：共同的尾结构一直就在题目里}
在过长的审计与检索之后，作者将讨论重新压回对象本身。他指出，本次 AF-by-discrete 比较应从共同的尾等价出发，离散部分表现为受控的尾类连接。需要跟踪的是 $G$ 与 $G_0$ 的持续边界点之间的共尾关系，以及相应回返在共同结构中的识别。

他的关键反问是：“你是不是搞复杂了？”这不是要求停止证明，而是要求停止把原题换成一个更宽泛的问题。若正在比较的持续边界代表本来就可以在同一 AF 结构中对应，那么首要动作应当是写出这个对应，而不是先撤去它，再设计恢复它的理论。

作者随后提出了一个更直接的结构解释：在这次固定的 AF-by-discrete 比较中，应先利用共同的尾结构与受控的离散拼接，把持续边界点放回相应的共尾类，再讨论两套标签下的边界运行与回返。这个解释直接针对此前审计忽略的题内结构，因此足以说明为什么原先把问题提升为一般 Cantor OE 是偏离。

但理由档案必须把“结构上有说服力”与“已在本文独立证明”分开。这里没有重新给出原论文中有限不可压缩性传递的完整形式证明，也没有把作者的陈述本身当作证明。两边标签可以不同，真正需要核对的是共同表示中的边界运行、共尾代表及回返约定；局部代表的统一字长界只属于固定有限生成集与所选局部表，不能无条件控制所有表示。

\claimstatus{作者指出共同 AF 尾结构的作用；讨论层面接受为结构解释，因为它直接恢复了此前被审计丢掉的题内结构。若把有限 IC 的传递作为正式数学结论引用，仍应回到 \cite[Section~4]{kua25} 的具体假设、定义与证明约定独立核验。这里不是一份新的独立证明。}

\subsection*{哪些复杂性本来不应成为主线}
此前，审计曾提议绕开有限 IC 的整体传递，转而尝试证明某种穿行计数的收缩，例如 $T_{n+d}(w)\leq qT_n(w)+C|w|$，其中 $q<1$。这里 $T_n(w)$ 表示相应层次的穿行计数。这可以是一类独立研究问题，却并不因为形式看起来更弱，就自动成为当前问题更好的路线。

在本案中，还没有先排除原有的直接几何解释，就把这种替代路线升成主角，实际增加了需要承担的工作。\textbf{寻找较弱目标只有在减少真实证明负担时才有价值；若它替换了本来可见的直接结构，就只是绕路。}

因此，本案撤回的不只是一个关于反例的判断，也包括一条不恰当的审计路线。不能在事后把所有绕行都包装成“方法论的必要胜利”。它曾经没有帮助作者解决问题，甚至使问题看起来比实际语境更危险。这一点应当被如实保存。

\subsection*{最后的纠正：方法论也会变成自己的模板}
作者随后指出，前面的许多审计是无用功。审计者承认，至少对这条有限不可压缩性传递的争论，主要错误来自将一般 Cantor OE 的困难带入已经具有额外结构的比较。

但仅仅为这个错误再增加一个名称，还不够。若每次方法失误都只产生更多条款，方法仍可能继续膨胀。作者因此给出了最后的约束：
\begin{casequote}
\textbf{“所以不要生搬方法论，大道至简。”}
\end{casequote}

这句话改变了 DCGI 的优先次序。方法不再以“完整执行审计”为目标，而以“是否需要审计”为前提。直接结构先行，额外分层随后。需要几分分辨率，就使用几分分辨率，不以最大分辨率作为默认美德。

同时，简单也不能由权威认同担保。作者的纠正给出了重新看对象的入口，入口之后仍然是数学。AI 从过度怀疑转成迅速认同，不应被误写成独立证明。真正应保存的，不是哪一方最后获胜，而是问题怎样从错误抽象回到可说明的结构。

\subsection*{本案确立与保留的结论}
第二次测试首先保留了一个正确纪律：轨道、胚、标记与词表达不应无说明地互换。但它同时限制了这条纪律的用法：必须先查看当前问题已经固定了哪些信息，不能为了强调信息损失而先丢弃这些信息。

它确立了语境漂移的识别方式，也使“结论强度匹配”同时约束赞成与反对。未见证明不是反例，作者澄清语境也不是取消所有证明义务。最重要的是，它把对象主权与至简原则提升为暂定的上位约束。

在具体研究层面，这场审计没有建立一个满足完整假设的反例，也不应被当作否定原论文传递主张的证据。讨论最终回到作者指出的共同 AF 尾结构、持续边界共尾性与回返对应这一结构解释。它在本判例中的作用，是纠正语境漂移并恢复原题的直接阅读路径；它在数学上的最终资格仍由原论文的确切约定与证明承担，而不是由这次对话中的认同承担。

\subsection*{判例登记：Orbit equivalence}
\begin{tabularx}{\textwidth}{@{}p{30mm}X@{}}
\toprule
触发事件 & 审计 $G_0$ finite IC 是否传递到有限生成的 $G\leq\FF(G_0)$，并讨论 orbit equivalence 在 AF-by-discrete 中到底保留什么。\\
原判断 & 由“一般 OE 太粗”迅速推向“finite IC transfer 可能失败”，随后又在作者纠正后过快转为认同。\\
偏离方向 & 一是 ambient-category drift：把具体 AF-by-discrete 语境换成任意 Cantor OE；二是把“尚未建立”说成“不一定成立”；三是作者纠正后把结构解释写得像独立证明。\\
纠正依据 & 反例资格要求迫使候选满足全部假设；作者指出共同 AF 尾结构与持续边界共尾性的题内作用，使此前被遗忘的结构重新进入论证。\\
影响条款 & 对象主权；结论资格；定位失败、最小修订；检索停止规则；至简原则。\\
当前状态 & 作者指出共同 AF 尾结构的作用，讨论接受其为\textbf{结构解释}；有限 IC 传递作为正式结论的独立验证仍应回到 \cite[Section~4]{kua25} 及其确切约定。本文不把对话中的认同登记为“已证”。\\
\bottomrule
\end{tabularx}


\section{两次测试为什么必须放在一起}
第一次测试告诉我们，不要让熟悉的实现占有概念。第二次告诉我们，不要让熟悉的方法占有对象。二者看起来方向相反，实际共同保护一种自由：使推理能够由对象的真实结构决定，而不是由先到场的模板决定。

只保留第一次经验，容易走向过度抽象：似乎每遇到一个对象，都应先释放它、推广它、另建一套机制。只保留第二次经验，又可能把“直接看结构”误解成无需追问依赖。两次测试一起，才使闭依与开生受到恰当约束。

\begin{center}
\small
\begin{tabularx}{\textwidth}{@{}p{25mm}XX@{}}
\toprule
 & \textbf{Just infiniteness} & \textbf{Orbit equivalence}\\
\midrule
起初的问题 & 性质能否离开树上的实现？ & 具体语境中的 IC 怎样传递？\\[5pt]
主要危险 & 把熟悉机制当作定义 & 把一般困难当作原题困难\\[5pt]
关键回归 & 回到定义与显式局部作用 & 回到共同 AF 结构与边界对应\\[5pt]
方法的变化 & 分离定义、机制与实现 & 对象主权与至简成为暂定的上位约束\\
\bottomrule
\end{tabularx}
\end{center}

因此，压力测试不能以发现多少问题计功，也不能以最终是否保住原结论计功。它应当留下更清楚的依赖、更准确的适用范围，或一个真实的新对象。如果留下的只有更长的解释，就应继续追问：这次测试究竟增加了什么？


\section{判例登记格式}
以后每次实际使用 DCGI，若事件足以修改或检验规则，先作为判例登记，不立即升格为新条款。固定记录六项：

\begin{tabularx}{\textwidth}{@{}p{30mm}X@{}}
\toprule
触发事件 & 什么具体问题、反例、误判或成功机制触发了这次登记？\\
原判断 & 在纠正发生以前，我们实际上怎样理解？\\
偏离方向 & 错误或不足发生在哪个方向，影响到原论证的哪一步？\\
纠正依据 & 是新的证明、满足全部假设的反例、文献结果、结构理由，还是仅仅作者陈述？\\
影响条款 & 哪些条款被新增、合并、削弱、确认或降级？\\
当前状态 & 已证、文献调用、候选、猜想、待闭合，或仅作为结构解释接受。\\
\bottomrule
\end{tabularx}

“当前状态”不能省略。它防止多年以后把一次有说服力的讨论认同误读成一次证明，也防止把尚未完成的候选路线因反复出现而自然升级为事实。


\chapter{基本条款}
\label{ch:articles}
条款不再作为一张平铺的原则清单出现。每一条末尾都标明它来自哪里、经历过怎样的压力测试。这样的 provenance 不是装饰，而是把“借来的理解必须可见”用在方法论自己身上。未经过独立判例的条款可以继续使用，但不得伪装成已经反复确认的规律。

\begingroup
\ctexset{section={beforeskip=18pt},subsection={beforeskip=12pt}}
\setlength{\parskip}{2pt plus 0.4pt}
\renewcommand{\lead}[1]{\par\noindent\textbf{#1}\par\nobreak\vspace{1pt}}

\section{问题怎样获得自己的形状}

\Article{问题与结构先于命名}
\lead{先让问题形成压力，再寻找能够承担它的语言。}
研究的起点未必是一句已经完整的命题。它可能是一个例子造成的不安，一处证明看起来过强的假设，或两个理论中反复出现的同一种作用。此时，最重要的不是迅速找到归类，而是保留这个问题尚未被现成语言替代的形状。

所谓“问题压力”，是对象持续要求解释的地方。它与一时的兴趣不同，也与对新术语的好奇不同。一个术语可以引起兴趣，但如果它没有回应任何实际困难，就不应仅凭吸引力取得研究中心的位置。

直觉可以在这一步先行。研究者可以先感到某种关系，而暂时说不清它。需要控制的不是直觉的出现，而是它取得结论资格的速度。\textbf{允许尚未命名的认识，不允许尚未验证的结论。}

命名因此有两个停止点：问题刚出现时，不以熟悉概念提前关闭它；一个 recurring role 尚未稳定时，也不急于给它制造专门术语。

一个反复出现的现象，应先被说明，再考虑命名。说明要回答它做什么、在什么条件下出现、与相邻概念怎样分开。若一句普通语言已经足够，没有必要再制造一个貌似专门的名词。

命名的收益应在后续推理中兑现。一个名称让多处论证共享同一个明确对象，就有价值。若它只是把一次性的描述包装成“结构”“原则”或“机制”，却没有减少任何歧义，那么它增加的是概念负担。

在正在形成的领域，这一点尤其重要。一个临时方便的说法可能很快被读者误认成标准定义。术语因此不只是写作风格，而是一项对未来使用者的责任。
\rulebasis{思想前史中的“问题先于概念”“延迟命名”；旧条款五、十一合并。}{思想前史推断为主，尚未经结构独立的数学判例专门检验。}

\Article{原生结构重建}
\lead{先重建“它是什么、怎样运作”，再决定采用哪一种解释。}
原生重建不要求从零开始，也不要求暂时遗忘已经学会的知识。它要求在接受一份外部说明之前，先尝试用自己能够检验的对象与关系恢复问题。

这种恢复可以很朴素：这里有哪些对象，哪些操作改变它们，什么保持不变，什么局部事实会产生整体后果？在数学中，这可能表现为一张关于子群、映射、商与作用的简图。在阅读其他理论时，也可以只是几句清楚的话。

“先识全貌”由此获得一个可操作的含义。全貌不是收集全部材料，而是先知道每一部分在整体中做什么。尚未理解的证明可以留白，但留白必须有位置。知道自己不知道哪一条依赖，往往比没有留白的复述更接近理解。
\rulebasis{思想前史；两次判例实际都先后要求回到对象、作用与依赖。}{两次判例均提供间接支持，但不是一条需要单独证明的数学命题。}

\Article{形式共鸣落在关系与作用上}
\lead{比较结构怎样组织结果，不只比较它们使用了哪些词。}
“两个领域都有边界”只能提出一个阅读入口，不能建立相容性。真正值得继续追踪的，是边界在两边是否承担可比较的作用：它是否区分局部与整体，是否限制延拓，是否改变后续的可能路径？

同样，“都有递归”不等于机制相同。有的递归用于复制局部作用，有的用于定义对象，有的用于得到缩短估计。名字相近，承担的依赖可能完全不同。

形式共鸣的价值，是让研究者在未完成细节之前先看到值得核验的关系安排。它不是证明的替代物，而是决定下一步读哪里、问什么的线索。正如听见一个主题的回归，不等于已经说明它为什么在此处改变了全曲的方向。
\rulebasis{思想前史的跨语言结构比较；判例一对 just infiniteness 机制的拆解。}{判例一直接支持“比较作用而非名词”；其他领域仍需后续判例复验。}

\section{外部知识怎样进入而不接管}

\Article{中心主权、先投影再吸收}
\lead{外部语言不能未经审查接管问题，自己的中心也不能免于修正。}
一个成熟理论往往同时提供对象、问题、范例与成功标准。吸收它时，研究者很容易连同这些标准一起接收，最后回答了一个漂亮的问题，却不是最初使自己不安的那个问题。

中心主权要求我们看见这种转换。某套语言解决了当前问题的一部分，可以借用。它提出了另一个更好的问题，也可以转向。但转向应当是明确的认识决定，不能藏在术语替换之中。

这个中心不是“我原先总是对的”。一旦反例击中中心区分，它就必须改变。所谓自主，不是只接受与自己相容的信息，而是知道自己为什么接受、为什么拒绝，以及在什么情况下必须重建。

中心主权决定谁拥有问题的组织权，先投影再吸收则说明外部理论怎样进入而不自动取得这项权力。

阅读外部理论时，可以先寻找它与当前结构图的对应位置：它是否补上了一条缺失的依赖，提供了一个此前没有的区分，或说明某个困难其实早已被解决？若找不到对应，也不必急于把它压进来。

吸收不是只有“全部接受”与“完全拒绝”两种选择。一个定理可以被直接调用，一种解释可以暂时保留，一个术语可以借用但重申范围，另一部分则可以与当前问题并置。不同部分拥有不同身份，正是清楚理解的表现。

这里的“投影”是阅读动作的比喻，不暗示已经构造了数学映射。它的要求很简单：\textbf{外部知识进入以后，应使问题中的某个位置更清楚，而不只是使叙述听起来更像某个学派。}
\rulebasis{思想前史的路径原创与语言适配；判例二的 ambient-category drift。}{判例二直接压力测试；与暂定“对象主权”宪法相连，仍等待结构不同的判例复验。}

\Article{借来的理解必须可见}
\lead{保存认识的来源，不把不同资格的内容压成一个“我知道”。}
自己从证明中重建出的机制、文献已经证明的定理、同行提供的解释，以及 AI 生成的候选猜想，不能只因为被放进同一段笔记就取得相同地位。

保持来源，不是为了争夺原创性，而是为了判断依赖的承重能力。一个来自文献的结果可以非常可靠，但它的使用条件需要确认。一个自己发现的关系可以非常有启发，却仍然需要证明。来源与真伪不是同一维度，但来源不清会使检验失去入口。

可以偶尔做一个简单测试：拿走眼前的解释，自己还能恢复到哪里？不能恢复的部分不必被羞于承认。它只说明那里仍有借用的支架，应当明确保留，而不是涂成已经独立掌握的墙面。
\rulebasis{思想前史的来源纪律；判例二中“作者纠正”与“独立证明”必须分开。}{判例二直接支持 provenance 的必要性。}

\Article{先区分，再统一}
\lead{相似提出问题，差异决定这种相似能走多远。}
面对两个看起来相容的对象，先找它们能够分开的情形。一个几乎满足条件却恰好失败的例子，常常比多个成功例子更能揭示真正的作用。

区分不是为了把一切拆得更碎。只有影响推理的差异才值得留下。若两种描述在当前目标下可以通过一个已知对应互换，就不必为了“精细”维持无效的分裂。

因此，先区分再统一，是一种可撤回的谨慎：先知道二者为何不同，再说明在哪一层可以把它们放在一起。它不是拒绝统一，而是使统一不再依靠含混。
\rulebasis{思想前史中 Similarity、Monogenesis、Compatibility 的分离；判例一的多机制阅读。}{思想前史与判例一共同支持；仍不是对所有比较问题的机械前置步骤。}

\section{依赖、对象与结论资格}

\Article{认知深度服从依赖}
\lead{优先深入那些一旦改变，就会改变整个结论的节点。}
学习不是把每个章节读到同样细。一个定义的细微约束可能决定全部证明，而一段很长的估计可能只是在实现已经清楚的界。阅读资源应当服从这种差别。

“关键依赖点”或 \emph{load-bearing node} 指的就是前一种位置。识别它的方法不是看作者用了多少强调词，而是问：把这一点拿走，结论还能沿原路成立吗？如果不能，这里值得深读。

这种策略并不允许永远略过技术。它只是安排次序：先恢复技术在整体中承担什么，再进入技术本身。困难段落因而不再是孤立的墙，而成为一座结构中可以定位的承重点。
\rulebasis{判例一对 load-bearing mechanism 的追问。}{判例一支持；判例二同时提醒“寻找承重点”不能演化成过度拆解。}

\Article{对象层级守恒}
\lead{推理可以改变对象，但不能不说明对象已经改变。}
词与群元素、群与群作用、轨道与胚群胚、有限近似与无限对象，都有各自的身份规则。一个词描述某个群元素，不代表它的每一个表达式具有相同的子词结构。两个作用拥有相同轨道，也不意味着它们给出了相同的局部实现。

这里的“守恒”不是禁止变换，而是要求交接清楚。若推理通过商、遗忘、限制、完备化或取极限，应说明什么被保留，什么被舍去，以及后面的结论需要的是哪部分。

不过，身份核对也必须服从对象主权。在一个已经固定了额外结构的比较中，不能先把额外结构忘掉，再宣布细节无从恢复。该保留的不能丢，该区分的不能混，二者是同一条纪律的两面。
\rulebasis{判例二早期关于 orbit、germ、word、action 的区分。}{判例二确认其必要性，同时限定其用法：不能先遗忘题目已经固定的额外结构。}

\Article{实现债务}
\lead{局部相容与整体实现之间若有距离，就明确承担跨过它的证明义务。}
几个局部模型分别可以实现，并不自动说明它们能在同一个整体中同时实现。每个有限阶段都有合适选择，也不自动说明这些选择彼此兼容。一个形式对应讲得通，也不自动说明它给出了实际机制。

“实现债务”，即 \emph{realization debt}，只是为这类具体义务保留一个可见的位置。它可能由一个延拓引理、紧致性论证、一致性条件或统一估计偿还。最重要的不是给债务命名，而是说出还缺什么。

若原定义或一个已经适用的结果本来就提供了整体实现，就应当直接调用，不必再虚设一道债务。这个条款用于标记真实的距离，不用于制造一条每次都要经过的桥。
\rulebasis{思想前史关于 compatibility 与 realization 的区分；判例一的非树实现问题。}{判例一支持其作为提醒工具；不得在已有整体实现时人为制造“债务”。}

\Article{结论资格与停止线}
\lead{论证说到哪里，结论就写到哪里。}
证明了局部实现，就不把它写成整体分类。计算了一个不变量，就不把它当作对象重建。找到了一个例子，就不称为普遍机制。说法的力度不能超过其实际支持。

同样，怀疑也有强度。没有看见证明，不等于证明不存在。现有论证尚未建立一个蕴含，不等于已经否定这个蕴含。只有反例、矛盾推导或相应结果，才足以承担否定的结论。

若一个强主张失败，可以保留仍被支持的弱主张，但必须明确撤回的是什么。不能在解释中退到较弱的位置，在摘要与结论中仍保留较强的声势。这既适用于作者，也适用于审计者。

“结论强度”与“停止线”回答的是同一问题：当前证据究竟允许我们把一句话写成什么身份。

合法结果可以是一条定理，一个构造，一个精确猜想，一个明确障碍，或一项可检验的研究方案。它也可以暂时停在一个尚未解决、但已被准确定位的问题上。

这里的 \emph{stop} 不是宣判问题不可解决，更不是因为问题长期存在就停止尝试。它是停止把当前证据提升到更强的资格。新的材料、新的机制或新的证明可以重新打开这条路径。

因此，“定理或止步”不要求所有认识都变成定理。它要求每一种认识使用自己的名字：猜想就是猜想，候选路线就是候选路线。认识的成熟，不只是能够向前走，也包括能够准确说明自己站在哪里。
\rulebasis{两次判例；尤其判例二中“没有证明”不能升级为“蕴含为假”。}{两次判例共同直接支持。执行层固定使用“已证 / 文献调用 / 候选 / 猜想 / 待闭合”等身份。}

\section{失败、检索与修订}

\Article{定位失败，最小修订}
\lead{先找出失败击中的位置，再决定修改的范围。}
若一条论证是从 $A$ 经 $B$、$C$ 到达 $D$，而问题发生在 $B$ 到 $C$ 的连接上，那么首先失效的是这条连接，不是四个对象本身。$C$ 可能有另一种证明，$D$ 可能只需要 $C$ 的较弱形式。

“失败的作用范围”是普通的依赖分析，不预设一个新的数学不变量。它可能落在定义、适用条件、局部推理、统一性或整体实现上，也可能根本落在审计者对题目的误读上。

局部化不是淡化问题。相反，准确指出一个错误影响哪些结论，比泛泛宣布“整个框架有问题”更有力量。它使修订能够进入对象，而不是停留在情绪与评价中。

只有先指出失败影响原论证的具体位置，才谈得上最小修订；否则“修订”很容易只是对模糊不安作出大范围反应。

一个反例出现以后，既不应立即推翻全部结构，也不应不停添加临时条件来保住原句。先决定哪个上游节点必须改变，再让改变沿依赖关系传播。

最小修订不是最小文字改动。有时只改一句话，却改变了整个定理的适用范围，那么下游必须一并更新。有时需要重写一段证明，却不影响结论，那么应该明确保留结论的身份。

修订应当留下简短记录：什么触发它，原来错在哪里，改动了什么，还有什么保持不变。这不要求建立庞大的管理系统。一段能够让未来的自己恢复判断的说明，通常已经足够。
\rulebasis{两次判例；其中“质疑与驳回都必须定位具体 proof step”的对称要求依附于暂定的对象主权。}{失败定位与最小修订由两次判例共同支持；对称 proof-step 地址要求直接来自判例二，当前随对象主权保持暂定。}

\Article{检索、探索与闭合各有任务}
\lead{在知道缺什么之后查找，而不是以不断阅读代替问题推进。}
检索可以由一个很具体的请求开始：某个蕴含是否已有证明，某个条件是否有标准名称，某个反例是否存在，或者一个性质是否有另一种实现。问题越清楚，文献越容易进入恰当的位置。

定向检索不排斥广泛阅读。漫游可以带来新的问题，但它与解决当前缺口是不同模式。若正在承担一个证明任务，就应知道下一次搜索将消除哪一种不确定性。

尤其要防止以搜索量冒充推进。连续命中同一段摘要，或在多个副本中重复找到同一个声明，并没有增加新的证明。检索的终点应当是能够返回当前问题，而不是生成下一轮更宽泛的检索。

检索是为某个结构缺口服务的动作；探索与闭合则决定当前允许怎样使用检索得到的材料。

探索时可以容纳类比、替代构造、尚不成熟的术语与相互竞争的解释。它们的价值在于扩大选择，而不是立即取得真理资格。

决定形成论文、证明或稳定说明以后，工作重心就改变了。定义与假设需要固定，引用与依赖需要准确，已知缺口需要关闭。此时不断新增词语、转换目标，可能只是在回避原来的证明义务。

两种模式可以切换，但切换应当可见。一个证明失败以后回到探索，并不是退步。把探索中的候选结果留在定理的位置上，才是角色混淆。
\rulebasis{思想前史；判例二中重复检索同一声明却不推进证明。}{判例二直接支持“检索停止”；探索/闭合的更一般区分仍主要来自思想前史与研究实践。}

\endgroup

\chapter{AI 执行模板}
\label{ch:templates}
可读版本在这里不再复制一套长操作手册。执行层由短模板派生，并且每一条要求都标明理由来源。AI 实际回答时不应朗读这些来源，也不应给步骤命名；模板只在后台约束协作。

\section{旧协议覆盖审计}
v0.2 的六份协议并非都能被最初拟定的四个执行模板自然覆盖。v0.3 先做覆盖审计，再决定迁移方式。初稿曾把协议 A“理解已有结果”判为由模板一完整覆盖；稳定化审计发现这一判断过高，因为理解一个结果与核验一份证明不是同一任务，因此新增模板六。

\begin{tabularx}{\textwidth}{@{}p{24mm}p{45mm}X@{}}
\toprule
旧协议 & v0.3 去向 & 覆盖判断\\
\midrule
A 理解已有结果 & 模板六“理解一个已有结果” + 模板一“证明与蕴含核验” & 理解层由模板六覆盖；若进一步核验证明，再转模板一。原“完整覆盖”判断撤回。\\
B 检验结构共鸣 & 新增模板五“结构共鸣核验” & 原四模板没有覆盖，因此新增，不假装已经被替代。\\
C 处理反例与失败 & 模板二“候选反例核验” + 模板一 & 完整覆盖。\\
D 准备推广 & 模板三“推广与替代机制” & 完整覆盖。\\
E 从闭依进入开生 & 模板三 + 本章“开生附注” & 只部分模板化；开生不能被强迫成固定程序。\\
F 与 AI 协作 & 六个模板的共同执行约束 + 模板四 & 分布式覆盖，不再独立成一份协议。\\
\bottomrule
\end{tabularx}

\section{所有模板共同遵守的七条协作规则}
以下七条是执行层的硬约束。它们不要求在回答里逐条复述。

\begin{enumerate}[leftmargin=2.2em,itemsep=5pt]
\item 先按题目给定的环境、结构和比较方式理解问题。不得自行换到更一般的语境；如果确实需要换，先说明理由。
\item 能直接论证的，就直接论证。不要引用方法论条款，不要给步骤命名。
\item 否定一个蕴含，必须给出满足全部假设的反例。否则只能说“尚未建立”。
\item 提出疑点时，要指出它作用在原证明的哪一步；驳回疑点时，也要指出使该步成立的结构理由。任何一侧指不出来，就不进入正式结论。
\item 作者纠正之后，要说明接受的依据：是结构上的理由，还是仅仅因为作者这样说了。不能因为被纠正就直接转为认同。
\item 每个结论都要标明身份：已证、文献调用、候选、猜想，还是待闭合。
\item 同一条检索结果出现第二次时，停止检索，改做推理。
\end{enumerate}

\section{模板一：证明与蕴含核验}
\begin{tabularx}{\textwidth}{@{}X p{53mm}@{}}
\toprule
\textbf{执行要求} & \textbf{理由来源}\\
\midrule
按题目原有的 ambient structure、对象与比较方式陈述目标，不擅自 generalize。 & 判例二；暂定对象主权。\\
先写最短直接论证；只有某一步确实卡住时才拆分。 & 判例二；暂定至简。\\
任何疑点必须写成“原证明的哪一个 implication / definition / estimate 需要什么，而当前只给了什么”；驳回疑点时也要指出使该步成立的具体结构。 & 判例二；依附于暂定的对象主权。\\
若作者纠正，明确记录“结构上接受”或“仅作者陈述，尚未独立验证”。 & 判例二的二次纠正。\\
输出结论时标明：已证 / 文献调用 / 候选 / 猜想 / 待闭合。 & 结论资格与停止线。\\
\bottomrule
\end{tabularx}

\section{模板二：候选反例核验}
\begin{tabularx}{\textwidth}{@{}X p{53mm}@{}}
\toprule
\textbf{执行要求} & \textbf{理由来源}\\
\midrule
先列原命题的全部 hypotheses，不允许只抓住其中一部分。 & 判例二。\\
候选反例须附核对表：逐项说明每个 hypothesis 是否满足，并给出依据。 & 判例二的碎片化候选；此前缺失 $G\leq\FF(G_0)$ 等关键资格时不能称为反例。\\
若有任何一项未核实，身份只能是“候选反例”；不得写成“蕴含不成立”。 & 判例二；结论资格。\\
真正否定蕴含时，必须明确展示“全部假设成立 + 结论失败”。 & 执行规则三。\\
若候选失败，记录它失败在哪个 hypothesis；不要把失败候选继续当作直觉证据。 & 定位失败、最小修订。\\
\bottomrule
\end{tabularx}

\section{模板三：推广与替代机制}
\begin{tabularx}{\textwidth}{@{}X p{53mm}@{}}
\toprule
\textbf{执行要求} & \textbf{理由来源}\\
\midrule
先分离目标性质、已知机制与历史 realization。 & 判例一。\\
若保持原机制，只核对原证明真正使用的数据，不把旧环境整套搬过去。 & 判例一；认知深度服从依赖。\\
若更换机制，明确说明新对象怎样承担原结论，而不是为旧术语寻找同义物。 & 判例一。\\
若整体实现尚未建立，明确登记 realization debt；已有直接实现时不得人为增加。 & 判例一；实现债务条款。\\
没有具体新机制时，可以停在候选/待闭合；不得为了“开生”强行 generalize。 & 判例一；结论资格与停止线。\\
\bottomrule
\end{tabularx}

\subsection*{开生附注}
旧协议 E 不能完全压成模板。Generative Independence 的关键恰恰是不让一个新模板重新规定未来必须怎样生成。执行时只保留两个检查：已经闭合的知识是否被分成“目标性质 / 机制 / realization”；新的方向是否拥有一个可检验的具体对象或障碍。除此之外，不规定必须产生新定理、新术语或新 framework。

\section{模板四：检索与结论资格}
\begin{tabularx}{\textwidth}{@{}X p{53mm}@{}}
\toprule
\textbf{执行要求} & \textbf{理由来源}\\
\midrule
每次检索必须对应一个明确缺口：找定义、找证明、找反例资格、找标准术语，或核对引用。 & 结构缺口牵引；判例二。\\
同一条结果第二次出现时，停止检索，转入推理；多个副本不算新的支持。 & 判例二的重复检索。\\
检索得到的“作者声称”与“正文已证明”必须分开登记。 & 判例二；借来的理解必须可见。\\
无法从文献或当前推理闭合时，输出“待闭合”，而不是用更多检索制造确定感。 & 结论资格与停止线。\\
\bottomrule
\end{tabularx}

\section{模板五：结构共鸣核验}
\begin{tabularx}{\textwidth}{@{}X p{53mm}@{}}
\toprule
\textbf{执行要求} & \textbf{理由来源}\\
\midrule
先让两边在各自原域中独立定义，不用一边的 vocabulary 反向切割另一边。 & 思想前史；路径原创与语言适配。\\
比较 relation / role / dependency，而不只比较相同名词或图像。 & 判例一；形式共鸣条款。\\
必须寻找至少一个 near-miss 或 separation condition，说明共鸣在哪一层停止。 & 先区分再统一。\\
若对应只换语言、不产生新的可检验区别，停止，不升级为结构同一或机制解释。 & 思想前史；结论资格。\\
结论身份明确标为结构相似、候选对应、已证关系或待闭合，不越级。 & 结论资格与停止线。\\
\bottomrule
\end{tabularx}


\section{模板六：理解一个已有结果}
这个模板回答“为什么这个结果在这些条件下成立”，而不是立即进入逐行 proof audit。它允许把成熟结果作为黑箱调用，但要求黑箱的作用可见。

\begin{tabularx}{\textwidth}{@{}X p{53mm}@{}}
\toprule
\textbf{执行要求} & \textbf{理由来源}\\
\midrule
先按题目给定的自然语境写清对象、主要 hypotheses 与 conclusion；不为追求一般性自行换环境。 & 闭依原则；暂定对象主权。\\
提炼一条\textbf{最短作用链}：哪些条件进入哪些机制或中间结果，最后怎样到达结论。只保留对理解有承重作用的节点。 & 闭依原则；认知深度服从依赖；暂定至简。\\
对每个暂作黑箱调用的结果，记录它的\textbf{输入与输出}，以及它在作用链中承担什么；不要求为了“理解”而把所有黑箱重新证明。 & 闭依原则；借来的理解必须可见。\\
对每个关键条件做去除测试：若去掉它，指出最短作用链在哪一步不再成立。若目前无法指出，只标“待闭合”，不得据此声称该条件必要或原结论为假。 & 判例一；判例二；结论资格与停止线。\\
区分“理解已经闭合”与“证明已经核验”。前者允许清楚标注的文献黑箱，后者才要求回到具体 statement、quantifier 与 proof step。 & 闭依原则；判例二。\\
最终把每个关键节点标成已证、文献调用、候选、猜想或待闭合，使读者能够看见理解在哪里结束、核验从哪里开始。 & 结论资格与停止线。\\
\bottomrule
\end{tabularx}

\chapter{自我修订、判例登记与版本史}
\label{ch:revision}
\section{理由档案的定位}
v0.3 明确区分“理由档案”与“执行规范”。理由档案保存一条规则为什么出现、经历过什么失败、当前处于什么资格；执行规范则尽可能短，只保留足以约束实际协作的模板。两者是单向派生关系：模板应从理由档案中产生，不能自行增加一套平行方法论。

这也改变了判例的位置。判例不再只是条款后面的例子，而是条款取得资格的主要来源之一。一个判例可以支持、限制或推翻规则；它不能因为被写进文档，就自动成为“成功案例”。

\section{暂定宪法怎样获得确认}
第三条“对象主权”和第四条“至简原则”在 v0.3 中降为暂定宪法，因为目前它们的直接压力测试来源主要是判例二。暂定条款仍约束当前执行，但不把一次有力案例伪装成永久充分依据。

若未来至少出现一个\textbf{结构不同}的判例，在没有复用 orbit-equivalence 同一机制的情况下再次独立要求同一上位原则，才考虑将其改为确认宪法。同一类型的问题重复十次，仍只能说明这条规则在同一结构族中稳定，不能用数量替代结构独立性。

\section{判例登记格式的调用}
固定的六栏判例登记格式已经列在第二编末尾“判例登记格式”一节，本编不再重复。今后的实际登记一律回到该格式；其中“当前状态”不可省略，“纠正依据”也必须区分证明、文献结果、结构解释与作者陈述，防止把事后的共识倒写成当时已经完成的证明。

\section{闭依的完成：能够说明，而非无穷追溯}
闭依不能被理解成对所有前提的无限追问。任何实际研究都在某个背景知识中进行，允许调用定义、基础结果与成熟工具。真正需要避免的，是关键处用不清楚的前提支撑清楚的结论。

对于当前任务，闭依完成意味着：我知道对象是什么，结论由什么机制产生，重要条件怎样进入，采用了哪些外部结果，以及一处改变会影响哪里。若某个依赖被明确作为假设接受，而不是伪装成已证明，它就有合法位置。

这使完成标准既严格又可达。它不要求一个人掌握整个领域，才有资格理解其中一条定理。它要求他能对自己的理解范围负责，并知道从哪里开始需要进一步展开。

\lead{闭依的标志：我知道这个结果为什么在这些条件下成立，而不只是知道它成立。}

\section{开生的完成：旧知识不再规定唯一形状}
开生并不以“与所有已知工作不同”作为标准。新的认识可以使用旧定理、旧技术与旧符号。它的独立性在于：这些材料的安排由当前问题决定，而不是由它们原先出现的顺序决定。

有时，开生表现为一个新构造。有时，它表现为发现原先被当作定义的条件其实只属于一种实现。还有时，它表现为一个准确的障碍，使研究不再重复无效路线。不同结果的强度不同，但都应当有可以指认的内容。

若尚未形成这样的内容，只获得一种形式共鸣，应当保留为线索。线索可以很重要，却不需要通过提前命名来抬高身份。一个问题被准确重写之后，已经为将来的证明保留了真实空间。

\lead{开生的标志：我能够使用这个理论，却不再必须按照它最初的实现方式安排新问题。}

\section{v0.3 相对于 v0.2 的变更记录}
本次升级遵守“Return 不是 Rewind”：v0.2 作为生成史在原工作档案中保留，不随本公开版分发；v0.3 记录为什么发生变化。

\begin{longtable}{@{}p{37mm}p{91mm}@{}}
\toprule
\textbf{变更} & \textbf{理由}\\
\midrule
\endfirsthead
\toprule
\textbf{变更} & \textbf{理由}\\
\midrule
\endhead
\bottomrule
\endfoot
可读版重新定位为“理由档案” & 它的主要价值已从“告诉 AI 怎么做”转为保存规则的来源、压力测试与资格。\\[4pt]
结构顺序改为“思想前史 → 宪法 → 判例 → 条款 → AI 模板 → 版本史” & 判例既然承担理由，就应先于由它们支持的条款出现。\\[4pt]
第三、四条宪法降为暂定 & 当前直接支撑主要来自判例二；需结构不同的独立判例再确认。\\[4pt]
基本条款由 17 条压为 12 条 & 合并“问题先于语言/机制先于术语”“中心主权/先投影再吸收”“结论强度/Theorem-or-Stop”“失败局部化/最小修订”“检索/探索与闭合”等重叠内容。\\[4pt]
每条条款增加“来源 + 压力测试状态” & 把 provenance 纪律用于方法论自己，防止预设规则伪装成已验证规则。\\[4pt]
v0.2 第三编旧六协议正文整体移除并退出规范层 & 协议不再作为可读版中的平行操作手册。先做覆盖审计，再由六个 AI 模板与开生附注承接；协议 A 因“理解/核验”分离新增模板六，协议 B 因原四模板未覆盖新增模板五，协议 E 仅部分模板化。\\[4pt]
删除 v0.2 附录 A 的“一页调用卡” & 可读版已经改为理由档案，执行摘要统一收束到 AI 模板；旧调用卡仍随 v0.2 历史版本在原工作档案中保留，不在 v0.3 正文重复，也不随本公开版分发。\\[4pt]
新增“质疑与驳回都必须定位具体 proof step” & 防止无地址的怀疑，也防止作者以身份替代结构回应。\\[4pt]
判例登记增加“当前状态” & 防止把讨论认同误读为证明，也防止候选路线在时间中自动升级。\\[4pt]
判例二重新降格具体数学结论 & 共同 AF 尾结构的作用作为作者提出并在讨论中接受的结构解释保存；正式 finite IC transfer 的资格仍回到原论文核验。\\[4pt]
减少纯文字判断的装饰性方框 & 理由档案以逻辑层级而不是视觉强调制造权重；保留少量真实引语的版式区分。\\
\end{longtable}

\standalonechapter{结语：让方法消失在清楚之中}
理解已有知识，需要一种向后的耐心。我们不能只带走结果，而把支撑它的条件、局部计算与实现劳动留在原地。创造新的知识，则需要一种向前的自由。我们不能因为已经知道一种成立方式，就让它规定所有可能的成立方式。

闭依与开生因此不是两种相反的性格，而是同一种认识责任的两个方向。向后，不轻率继承。向前，不机械复制。两者之间，没有一个可以替人完成判断的总框架，只有对象不断要求我们作出的具体区分。

第一次压力测试让一个概念从树的影子里走出来。它没有否定树的价值，而是使树上的工作成为可以重新理解、重新安排的材料。第二次压力测试则让方法从自己的影子里走出来。它没有否定审计的价值，而是使审计重新服从其原本应当帮助的对象。

于是，这套体系最后留下两条已经构成定义轴线的承诺——对已有依赖负责，不让已有依赖占有未来——以及两条正在接受进一步判例检验的上位约束：尊重对象已经给出的结构，只使用当前问题所需的分辨率。

当这些承诺真正进入工作，方法论就不必始终站在舞台中央。它应当在必要时提醒，在混淆处区分，在失败后帮助修订，然后退开，让读者看见那个已经变得清楚的对象。

\begin{casequote}
理解，是对已有依赖负责。\par
创造，是不让已有依赖占有未来。\par
方法论服从对象。大道至简。
\end{casequote}

\noindent 成熟的 DCGI，最终应当消失在对象本身的清晰结构中。

\appendix
\ctexset{chapter={name={附录,},number=\Alph{chapter}}}
\ctexset{chapter={format=\LARGE\sffamily\bfseries\color{ink},nameformat=\Large\sffamily\color{muted},aftername=\quad,beforeskip=0pt,afterskip=18pt}}
\chapter{术语与阅读定位}
\label{app:terms}
这里的英文词用于与研究讨论中的说法对应，不另外构成一套必须掌握的专门语言。除已有数学术语外，其余词汇都是认识活动的描述。

\small
\begin{longtable}{@{}p{43mm}p{85mm}@{}}
\toprule
\textbf{用语} & \textbf{在本文中的意思}\\
\midrule
\endfirsthead
\toprule
\textbf{用语} & \textbf{在本文中的意思}\\
\midrule
\endhead
\bottomrule
\endfoot
Dependency Closure & 闭依。使当前问题的重要依赖有明确来源与适用条件，不要求无穷追溯全部基础。\\[4pt]
Generative Independence & 开生。使已有知识成为可重新安排的材料，不把既有实现当作未来的唯一模板。\\[4pt]
Problem pressure & 问题压力。对象中持续要求解释的困难、差异或尚未表述的关系。\\[4pt]
Native reconstruction & 原生结构重建。先用自己可以检查的对象与关系恢复问题，不等于拒绝外部知识。\\[4pt]
Projection before absorption & 先投影，再吸收。先辨认外部理论回应当前问题的什么位置。这里不是数学投影。\\[4pt]
Borrowed understanding & 借来的理解。尚需依靠特定外部解释或结果的认识，其来源应当可见。\\[4pt]
Load-bearing node & 关键依赖点。改变它会改变结论或主要论证的位置。\\[4pt]
Object-level conservation & 对象层级守恒。变换、比较与识别必须说明所处理的对象及其身份规则。\\[4pt]
Realization debt & 实现债务。局部相容走向整体实现时，确实尚未完成的证明义务。\\[4pt]
Failure support & 失败的作用范围。一个错误实际击中的依赖及其下游影响，不是额外定义的数学不变量。\\[4pt]
Near-miss & 近失例。大部分结构相似，却因关键条件不成立而与成功情形分离的例子。\\[4pt]
Ambient-category drift & 语境漂移。未经说明改变研究环境，再把新环境中的困难记在原问题名下。\\[4pt]
Exploration / closure & 探索与闭合。前者提出可能，后者承担已经提出的主张。\\[4pt]
Conclusion status / stop line & 结论资格与停止线。用“已证、文献调用、候选、猜想、待闭合”等身份说明当前到达哪里，不把停止升级主张等同于放弃问题。\\
\end{longtable}
\normalsize

\standalonechapter{资料与版本说明}
v0.3 将本文明确定位为 DCGI 的理由档案。它保存思想前史、宪法资格、两次数学压力测试、由判例支持的基本条款、AI 执行模板的来源，以及版本如何修订。执行模板来自理由档案，不另立一套与正文平行演化的方法。

数学判例继续保留不同结论的身份。Just infiniteness 一案中的文献结果、文内小证明与 AF-by-discrete 候选方向分开登记。Orbit equivalence 一案尤其不把作者纠正后的讨论认同写成独立证明：共同 AF 尾结构的作用被记录为具有明确结构理由的解释，正式 finite IC transfer 仍由原论文的具体假设与证明约定承担。

思想前史仍取材于彭国欣与 AI“钟上的降雪”围绕路径原创、问题先于概念、语言适配、Monogenesis、Return、Interface 与去中心化整合所写的一组文章。它们说明问题怎样生成，不为数学判例背书。

v0.2 作为上一阶段文本在原工作档案中保留，不随本公开版分发。v0.3 的详细变更见第五编“v0.3 相对于 v0.2 的变更记录”。

下列文献用于数学背景、术语核对与思想前史定位。它们不为 DCGI 的全部认识论主张提供外部权威担保。

\clearpage
\phantomsection
\addcontentsline{toc}{chapter}{参考文献与阅读线索}
\markboth{参考文献}{参考文献}
\begin{thebibliography}{9}
\small
\setlength{\itemsep}{9pt}
\bibitem{bgs03}
Laurent Bartholdi, Rostislav I. Grigorchuk, Zoran Šunić.
\newblock \emph{Branch Groups}.
\newblock In: \emph{Handbook of Algebra}, Vol.~3, North-Holland, 2003, pp.~989--1112.
\newblock 公开版本：\href{https://arxiv.org/abs/math/0510294v2}{arXiv:math/0510294v2}。
\newblock 阅读线索：Introduction 与 Chapter~5，尤其 Definition~5.1、Theorem~5.2、Lemma~5.3 及 Theorem~5.6。这里也是本文核对 Wilson 与 Grigorchuk 相关结构理论背景的入口。

\bibitem{nek18}
Volodymyr Nekrashevych.
\newblock Palindromic subshifts and simple periodic groups of intermediate growth.
\newblock \emph{Annals of Mathematics}, 187(3), 2018, pp.~667--719.
\newblock \href{https://annals.math.princeton.edu/2018/187-3/p02}{期刊页面}。
\newblock 阅读线索：从 Cantor 动力系统及碎片化构造理解单纯群路线，不将它与一般分支机制混同。

\bibitem{kua24}
Zheng Kuang.
\newblock Growth of groups with incompressible elements, I.
\newblock \emph{Groups, Geometry, and Dynamics}, 2026, published online first.
\newblock \href{https://doi.org/10.4171/GGD/954}{doi:10.4171/GGD/954}。
\newblock 公开预印本：\href{https://arxiv.org/abs/2402.16238v3}{arXiv:2402.16238v3}。
\newblock 本书关于 reduced word 与两种估计的定义编号依上述预印本 v3 的 Section~4，特别是 Definition~4.2 与 Remark~4.5，避免不同版本编号混用。

\bibitem{kua25}
Zheng Kuang.
\newblock Growth of finitely generated subgroups of the topological full groups of inverse semigroups of bounded type.
\newblock \href{https://arxiv.org/abs/2409.19491v2}{arXiv:2409.19491v2}, 29 May 2025.
\newblock 阅读线索：Introduction；Definitions~2.17--2.18；Section~4 中的 Theorem~4.1、Corollary~4.3 与 Definition~4.3。本文记录的有限 IC 争论以这一版本为讨论对象。


\bibitem{peng-problem}
彭国欣、AI“钟上的降雪”.
\newblock 《问题先于概念：AI时代为何要保护自己的理解路径？——从路径原创到“延迟命名”的思考》.
\newblock 2026 年 9 月 27 日.
\newblock 阅读线索：问题先于命名、概念的压缩能力，以及 AI 时代“延迟命名”为什么成为主动的认识纪律。

\bibitem{peng-path}
彭国欣、AI“钟上的降雪”.
\newblock 《理解有路，真理无路——从“路径原创”到克里希那穆提的“无路之国”：存在几何学的一次结构推导》.
\newblock 2026 年 9 月 27 日.
\newblock 阅读线索：理解必然有路径、路径只能在理解中生成，以及既有道路为何不能取得对未来理解的先验所有权。

\bibitem{peng-language}
彭国欣、AI“钟上的降雪”.
\newblock 《语言之间不是等级关系，而是问题—语言之间存在适配关系》.
\newblock 2026 年 9 月.
\newblock 阅读线索：描述语言通过选择差异而工作；问题—语言适配比寻找“最高语言”更重要；模型不能因此取得对象的垄断解释权。

\bibitem{peng-detour}
彭国欣、AI“钟上的降雪”.
\newblock 《为什么我们必须绕远路？——从误解、投影到 Monogenesis 与 Return》.
\newblock 2026 年 9 月 15 日.
\newblock 阅读线索：Generation--Differentiation--Projection--Path--Return；尤其“Return 不是 Rewind”，返回必须把差异与路径带回来。

\bibitem{peng-return-boundary}
彭国欣、AI“钟上的降雪”.
\newblock 《Return 的边界在哪里？——关于可返回性的追问：从 Monogenesis 到 Pushout 与新的共同结构》.
\newblock 2026 年 9 月 16 日.
\newblock 阅读线索：共同来源不自动等于可返回；Similarity、Monogenesis 与 Compatibility 的分离；新 interface 是否可以重新创造共同结构。

\bibitem{peng-center}
彭国欣、AI“钟上的降雪”.
\newblock 《当中心消失之后，结构如何重新整合？——从三阶意识、Monogenesis 到 AF-by-discrete 的 Interface、Compatibility 与生成性 Return》.
\newblock 2026 年 9 月 17 日.
\newblock 阅读线索：去中心化之后怎样避免碎片化；统一可以来自共同组织条件，而不必恢复一个最终中心。

\bibitem{peng-babel}
彭国欣、AI“钟上的降雪”.
\newblock 《AI可否克服语言分裂，重建人类巴别塔？——视角分裂、共同中心与整体化中的意识场》.
\newblock 2026 年 9 月.
\newblock 阅读线索：AI 可以重接语言接口，却不能消除观察位置，也不能由技术性可互译直接取得问题主权。

\end{thebibliography}
\end{document}
